\documentclass[12pt,oneside,a4paper]{report}
\newcommand{\annexname}{Anexo}
\makeatletter
\newcommand\annex{\par
 \setcounter{chapter}{0}%
 \setcounter{section}{0}%
 \gdef\@chapapp{\annexname}%
 \gdef\thechapter{\@Alph\c@chapter}}
\makeatother
\usepackage[utf8]{inputenc}
\usepackage[spanish, es-tabla, es-lcroman]{babel}
\usepackage{hyperref}

\usepackage[labelsep=period]{caption}
\usepackage{amsmath}
\usepackage{amsfonts}
\usepackage{graphicx}
\usepackage{amssymb}
\usepackage{gensymb}
\usepackage{changepage}
\usepackage{times}
\usepackage{setspace}
%\usepackage[capposition=top]{floatrow}
\usepackage[version=3]{mhchem}
\usepackage{textgreek}
\usepackage{geometry}
\usepackage{appendix}
\usepackage{sectsty}
\usepackage{lipsum}
\usepackage{fancyhdr}
\setlength{\headheight}{15.2pt}
\pagestyle{fancy}

\renewcommand{\appendixname}{Anexos}
\renewcommand{\appendixtocname}{Anexos}
\renewcommand{\appendixpagename}{Anexos}



\renewcommand{\baselinestretch}{1.5} 

\setcounter{secnumdepth}{3}

\geometry{a4paper, left=40mm, right=25mm,top=35mm,bottom=25mm,headsep=20mm}



\allsectionsfont{\normalsize}
\chapternumberfont{\centering \large}


\lhead{}
\chead{} 
\rhead{\thepage} 
\cfoot{}
\rfoot{Nombre del tesista}
\renewcommand{\footrulewidth}{0.5pt}

\begin{document}







\thispagestyle{empty}
\pagenumbering{gobble}
\setstretch{1.0}
\begin{center}

\Large\textbf{{TITULO DE LA TESIS\\}}

\vspace{1cm}

\includegraphics[scale=0.65]{logo_una}\\

\vspace{1cm}

\large{\textbf{
NOMBRE DEL TESISTA\\
}}
\vspace{1cm}

TESIS\\

\vspace{1cm}

Presentada a la Facultad de Ciencias Químicas de la Universidad Nacional de Asunción como requisito para la obtención del grado de INGENIERO QUÍMICO

\vspace{2cm}

San Lorenzo - Paraguay\\
AÑO

\vspace{1cm}






\pagebreak
\thispagestyle{empty}

\Large\textbf{{TITULO DE LA TESIS\\}}

\vspace{4cm}

\large{\textbf{
NOMBRE DEL TESISTA\\
}}

\vspace{4cm}
Orientador: Ing. Quím. \textbf{NOMBRE DEL ORIENTADOR}, M.Sc.\\

Co-orientador: Ing. Quím. \textbf{NOMBRE DEL CO-ORIENTADOR}, M.Sc.

\vspace{3cm}
Tesis presentada a la Facultad de Ciencias Químicas de la Universidad Nacional de Asunción como requisito para la obtención del grado de INGENIERO QUÍMICO

\vfill
San Lorenzo - Paraguay

AÑO

\newpage
\thispagestyle{plain}

\Large\textbf{{TITULO DE LA TESIS\\}}

\vspace{3cm}

\large{\textbf{
NOMBRE DEL TESISTA
}}\\
\vspace{3cm}
\end{center}

Aprobado en fecha X de XXXXX de 20XX

\vspace{2cm}

Tribunal examinador:

\begin{itemize}

\item Ing. Quím. Miembro de mesa 1, PhD. \hfill FCQ/UNA

\item Ing. Quím. Miembro de mesa 2, PhD. \hfill FCQ/UNA
\item Ing. Quím. Miembro de mesa 3, Msc. \hfill FCQ/UNA
\end{itemize}
\begin{center}
\vfill


.....................................................................

Ing. Quím. NOMBRE DEL ORIENTADOR, MSc.


Orientador

\end{center}

\newpage
\thispagestyle{plain}
\pagenumbering{roman}
\setcounter{page}{3}
\rhead{}
\chead{}
\lhead{\thepage}
\rfoot{}
\cfoot{}
\lfoot{}
\vspace{2cm}
\hfill \textit{Dedicatoria}\\
\vspace{1cm}


\begin{flushright}
\textit{Aquí va la dedicatoria}
\end{flushright}

\newpage
\thispagestyle{plain}
\begin{center}
\textbf{AGRADECIMIENTOS}\\
\vspace{1cm}

Primer agradecimiento.\\

\vspace{1cm}

Segundo agradecimiento\\

\vspace{1cm}

Tercer agradecimiento\\

\vspace{1cm}

Cuarto agradecimiento\\


\newpage
\thispagestyle{plain}
\textbf{AGRADECIMIENTOS INSTITUCIONALES}\\

\vspace{1cm}

A la Facultad de Ciencias Químicas de la Universidad Nacional de Asunción, por el uso de sus instalaciones y equipos.

\vspace{1cm}

Segundo agradecimiento institucional

\vspace{1cm}

Tercer agradecimiento institucional

\vspace{1cm}


\end{center}



%\frontmatter

\newpage
\thispagestyle{plain}

\begin{center}
\large{\textbf{{TITULO DE LA TESIS}}}\\
\end{center}
\setstretch{1.5}

\normalsize
\
 \hfill Autor: TESISTA

\hfill Orientador: NOMBRE DEL ORIENTADOR

\hfill Co-Orientador: NOMBRE DEL CO-ORIENTADOR


\textbf{Resumen}


\noindent 

\textit{Palabras clave: palabras, clave.}

\newpage
\thispagestyle{plain}
\setstretch{1.0}
\begin{center}
\large{\textbf{TITULO DE LA TESIS EN INGLES}}\\
\end{center}
\setstretch{1.5}

\normalsize{

\hfill Author: NOMBRE DEL TESISTA

\hfill Advisor: NOMBRE DEL ORIENTADOR

\hfill Co-Advisor: NOMBRE DEL CO-ORIENTADOR

\textbf{Summary}




\textit{Keywords: palabras, clave}

\newpage

\lhead{}
\rhead{\thepage}

\tableofcontents

\pagebreak

\listoffigures

\pagebreak

\listoftables













\pagebreak
\lhead{Capítulo \thechapter}
\chead{} 
\rhead{\thepage} 
\cfoot{}
\rfoot{Nombre del tesista}
\renewcommand{\footrulewidth}{0.5pt}


\pagenumbering{arabic}
%\mainmatter
\chapter{Introducci\'on}





 
\chapter{Revisi\'on de la literatura}


\chapter{Metodolog\'ia}


\chapter{Resutados y discusión}


 
 \chapter{Conclusiones}






\annex
\addcontentsline{toc}{chapter}{Anexos}
\lhead{Anexo \thechapter}

\chapter{Título del anexo 1}


\chapter{Título del anexo 2}


\newpage
\setstretch{1.0}
\lhead{Biliografía}
\begin{thebibliography}{9}

\bibitem Aquí va la bibliografía

\end{thebibliography}

\appendix


\chapter{Título del apéndice 1}
\pagenumbering{Alph}




\end{document}